% s-042-p042-m002; original PDF 42, 43
\Aphorism{s-042-p042-m002}{29}
Vielleicht ist dies das sicherste Zeichen der vornehmen Seele, dass sie sich nicht auf Handel und Vortheil versteht, im Grossen wie im Kleinen. Nicht als fehlte ihr dazu das nöthige \Foreign{ingenium}; wo es sein muss, führt der höhere Mensch eine anspruchsvolle Situation (nur nicht gerade eine Intrigue) mit ebensoviel Feinheit und Ueberblick durch wie nur irgend ein Practicus; im Gegentheil scheint es, als ob gerade die inferioren Leute sich nicht einmal im äusseren Leben zweckmässig einzurichten wüssten, wo sie doch eigentlich zu Hause sind. Nicht der Geist, das Handwerkszeug, sondern Ernst und Interesse geht dem vornehmen Menschen in utilitarischen Angelegenheiten ab und haftet dem Alltagsmenschen an; hierbei nicht ausgeschlossen, dass die ungern zugreifende, behandschuhte Hand an der richtigen Stelle zufasst, die von gelenken Allerweltsfingern verfehlt wird. — Die vornehme Seele lässt sich übervortheilen, sie findet es nicht unbillig, ihre Erfahrungen theurer einzukaufen als Jedermann. Sie ist fern von der „plumpen Sich-Recht-geberei“, die jeden ihr Interesse kreuzenden Willen böse nennt; sie ist weniger fern davon, ihre Feinde zu \Sperr{lieben}. Sie betet Fehlgriff und Enttäuschung an; hinter einem Verluste herzuschimpfen würde ihr erst der eigentliche Verlust sein, nämlich der der guten Manieren. Sie will sich nicht bewahren, sie sucht nicht nach Gründen und Rechten, sich zu bewahren, sie frägt nicht, ob an ihrer Stelle Andere sich zu opfern hätten und ob ihr Opfer Frucht tragen werde. Man achte darauf, wenn ein Krieg ausbricht, wer sich für die nicht versäumte Pflicht am theuersten bezahlen lässt, mit einem Sack voll Kupfergeld und moralischer Scheidemünze, man achte darauf, wer am ehesten der Versuchung anheimfällt und erliegt, sich seiner Pflicht zu entziehen: gewiss nicht Diejenigen, an deren Fortbestand etwas gelegen ist, die ein „Recht“ hätten sich fernzuhalten, die vornehmen und seltenen Menschen: — diese stehen und fallen in vorderster Linie! Sondern die Ueberflüssigen, die Viel-zu-Vielen, die nur als „Kanonenfutter“ dasein sollten, weil sie anders dem Vaterlande ihre Schuld nicht bezahlen können: diese wissen vor der Schlacht oder in der Schlacht davonzukommen. Es ist gegen den Krieg zu sagen, dass er in der biologischen wie in der Cultur-Auslese einen Schritt rückwärts bedeutet: wie er die tüchtigsten Exemplare opfert, um die Schwachen und Gebrechlichen zu retten, so verhilft er den edlen, leidenschaftlichen Naturen, die sich gerne verschwenden, zum gesuchten Untergange und den engherzigen, selbstsüchtigen Feiglingen zum Ueberleben.
\par

